\documentclass[11pt,a4paper]{article}
\usepackage{amsmath}
\usepackage{amssymb}
\usepackage{graphicx}
\usepackage{fontspec}
\usepackage{babel}
\usepackage[margin=2.5cm]{geometry}
\usepackage{enumitem}
\usepackage{booktabs}
\usepackage{xcolor}
\usepackage{hyperref}
\usepackage{abstract}

\definecolor{coral}{HTML}{cc785c}
\definecolor{ink}{HTML}{141413}
\definecolor{muted}{HTML}{6c6a64}

\hypersetup{colorlinks=true, linkcolor=coral, urlcolor=coral, citecolor=coral}

\babelprovide[import, onchar=ids fonts]{thai}
\babelprovide[import, onchar=ids fonts]{english}

\babelfont{rm}{Noto Sans}
\babelfont[thai]{rm}{Noto Sans Thai}
\babelfont{tt}{Noto Sans Mono}
\babelfont[thai]{tt}{Noto Sans Thai}

\directlua{
  local glyph_id = node.id("glyph")
  local penalty_id = node.id("penalty")

  local function is_thai(c)
    return c >= 0x0E01 and c <= 0x0E5B
  end

  local function is_thai_combining(c)
    return (c >= 0x0E31 and c <= 0x0E3A) or (c >= 0x0E47 and c <= 0x0E4E)
  end

  luatexbase.add_to_callback("pre_linebreak_filter", function(head)
    for n in node.traverse_id(glyph_id, head) do
      local prev = n.prev
      if prev and prev.id == glyph_id and is_thai(prev.char)
         and is_thai(n.char) and not is_thai_combining(n.char) then
        local p = node.new(penalty_id)
        p.penalty = 0
        node.insert_before(head, n, p)
      end
    end
    return head
  end, "thai_break")
}

\emergencystretch=1em

%% ─── Title metadata ───────────────────────────────────────────
\title{\textbf{ชื่อบทความ}\\[6pt]
  {\large Subtitle in English}}
\author{ชื่อผู้เขียน\\
  {\small สถาบัน · อีเมล}}
\date{\today}

\begin{document}

\maketitle

\begin{abstract}
บทคัดย่อภาษาไทย (150--200 คำ) อธิบายปัญหา วิธีการ และผลลัพธ์หลัก
ของงานวิจัยนี้อย่างกระชับ
\end{abstract}

\noindent\rule{\textwidth}{0.5pt}
\vspace{8pt}

%% ─── 1. Introduction ──────────────────────────────────────────
\section{บทนำ}

อธิบายที่มาและความสำคัญของปัญหา บอกว่างานนี้แก้ปัญหาอะไร
และมีส่วนสนับสนุนต่อองค์ความรู้อย่างไร

\subsection{แรงจูงใจ}

อธิบายว่าทำไมปัญหานี้จึงสำคัญ

\subsection{วัตถุประสงค์}

\begin{enumerate}[nosep]
  \item วัตถุประสงค์ที่ 1
  \item วัตถุประสงค์ที่ 2
  \item วัตถุประสงค์ที่ 3
\end{enumerate}

%% ─── 2. Related Work ──────────────────────────────────────────
\section{งานวิจัยที่เกี่ยวข้อง}

สรุปงานวิจัยก่อนหน้าที่เกี่ยวข้อง จัดกลุ่มตามแนวทาง
และชี้ให้เห็นช่องว่างที่งานนี้เติมเต็ม

%% ─── 3. Methodology ───────────────────────────────────────────
\section{วิธีการ}

\subsection{ภาพรวมระบบ}

อธิบายสถาปัตยกรรมหรือวิธีการโดยรวม

\subsection{รายละเอียดวิธีการ}

อธิบายขั้นตอนหรือองค์ประกอบแต่ละส่วน

%% ─── 4. Experiments ───────────────────────────────────────────
\section{การทดลองและผลลัพธ์}

\subsection{การตั้งค่าการทดลอง}

\begin{itemize}[nosep]
  \item Dataset: ชื่อและขนาดชุดข้อมูล
  \item Baseline: วิธีการเปรียบเทียบ
  \item Metric: ตัวชี้วัด
\end{itemize}

\subsection{ผลลัพธ์หลัก}

\vspace{8pt}
\noindent
\begin{tabular}{@{}llll@{}}
  \toprule
  \textbf{Method} & \textbf{Metric 1} & \textbf{Metric 2} & \textbf{Metric 3} \\
  \midrule
  Baseline A   & 0.00 & 0.00 & 0.00 \\
  Baseline B   & 0.00 & 0.00 & 0.00 \\
  Ours         & \textbf{0.00} & \textbf{0.00} & \textbf{0.00} \\
  \bottomrule
\end{tabular}
\vspace{8pt}

%% ─── 5. Discussion ────────────────────────────────────────────
\section{อภิปราย}

วิเคราะห์ผลลัพธ์ อธิบายสาเหตุที่วิธีของเราดีกว่า
และชี้ข้อจำกัดหรือกรณีที่ล้มเหลว

%% ─── 6. Conclusion ────────────────────────────────────────────
\section{สรุป}

สรุปผลสำคัญและทิศทางวิจัยในอนาคต

\end{document}
